\documentclass[11pt,a4paper]{article}
\usepackage[margin=2.5cm]{geometry}
\usepackage{amsmath,amssymb}
\usepackage{booktabs}
\usepackage{parskip}
\usepackage{hyperref}
\usepackage{xcolor}
\usepackage{graphicx}
\usepackage{siunitx}

\hypersetup{colorlinks=true,linkcolor=blue,citecolor=blue,urlcolor=blue}

\newcommand{\gtrue}{\gamma_{\rm true}}
\newcommand{\ghat}{\hat\gamma}
\newcommand{\dm}{\Delta m}
\newcommand{\dmref}{\Delta m_{\rm ref}}
\newcommand{\Mpc}{\,\mathrm{Mpc}}
\newcommand{\Msun}{\,M_{\odot}}

\title{\textbf{Diagnosis of the $\ghat \approx -0.5$ Plateau in the Gamma Scan}\\[0.4em]
\large Destacking project --- model-misspecification analysis}
\author{destacking project}
\date{2026 May 26}

\begin{document}
\maketitle

\begin{abstract}
The 11-point post-hoc-bootstrap gamma scan
(\texttt{gamma\_scan\_outputs/run\_2026-05-24T10-19-32/}) recovers
$\ghat \approx -0.5$ for every $\gtrue \le -0.6$, producing a plateau in the
transfer function $\ghat(\gtrue)$. We show that the bias lives in the likelihood
itself, not the sampler, and is not a parametrisation artefact: even when the
inference uses the \emph{exact per-cluster truth concentrations}, the spherical
NFW model is beaten by parametric $\gamma = -0.5$ by 3851 log-likelihood units
at $\gtrue = -1$. The root cause is the mismatch between the
\emph{prolate-dominated} triaxial halo population used in the simulation
(47\% prolate, 31\% triaxial, 21\% oblate; median triaxiality $T = 0.643$) and
the spherical NFW model used for inference. Negative $\gtrue$ pushes the
implied $c$-distribution into bimodal tails where the model misspecification
is worst; the plateau at $\gamma \approx -0.5$ is the value at which the
$c$-distribution stays moderate enough that orientation-averaged triaxial
projections look approximately like spherical NFW. Because the real
observational data carry no information on per-cluster triaxiality, this
bias cannot be removed within the project scope, and $\ghat$ must be
interpreted as an \emph{effective}, projection-averaged slope, not the
intrinsic $\gtrue$.
\end{abstract}

\tableofcontents
\newpage

% ============================================================
\section{The phenomenon}
% ============================================================

The reference run scans $\gtrue \in \{-1.0,-0.8,\ldots,+1.0\}$ (11 values, 10000
clusters, 50 post-hoc bootstrap draws per $\gtrue$).  The ensemble medians
$\ghat = \mathrm{median}_b(\ghat^{(b)})$ define the transfer function
shown in Table~\ref{tab:scan}.  For $|\gtrue| \gtrsim 0.5$ the estimator
asymptotes onto a one-sided floor near $\ghat \approx \pm 0.5$.

\begin{table}[htbp]
  \centering
  \begin{tabular}{rrrr}
    \toprule
    $\gtrue$ & $\ghat$ & $\sigma_{\rm boot}$ & bias\\
    \midrule
    $-1.0$ & $-0.548$ & $0.171$ & $+0.452$\\
    $-0.8$ & $-0.542$ & $0.122$ & $+0.258$\\
    $-0.6$ & $-0.410$ & $0.127$ & $+0.191$\\
    $-0.4$ & $-0.494$ & $0.057$ & $-0.094$\\
    $-0.2$ & $-0.516$ & $0.242$ & $-0.316$\\
    $\phantom{-}0.0$ & $-0.165$ & $0.342$ & $-0.165$\\
    $+0.2$ & $+0.453$ & $0.340$ & $+0.253$\\
    $+0.4$ & $+0.174$ & $0.268$ & $-0.226$\\
    $+0.6$ & $+1.339$ & $0.618$ & $+0.739$\\
    $+0.8$ & $+0.560$ & $0.274$ & $-0.240$\\
    $+1.0$ & $+0.750$ & $0.251$ & $-0.250$\\
    \bottomrule
  \end{tabular}
  \caption{Transfer function from the 11-point bootstrap scan. The bias does
  not vanish at large $|\gtrue|$; instead $\ghat$ saturates near $\pm 0.5$.}
  \label{tab:scan}
\end{table}

\noindent This report concerns the negative branch; the positive side is
symmetric by construction of the $c(\dm)$ parametrisation.

% ============================================================
\section{The likelihood itself is biased --- not the sampler}
% ============================================================

The first thing to rule out is a nautilus pathology.  We hold
$(c_0,\alpha,\beta,\log M_{\rm norm})$ fixed at the OLS-truth values
(the coefficients that exactly reproduce the simulated
$\log_{10}(c_{\rm truth})$ regression) and scan $\gamma$ over a fine grid using
the same \texttt{log\_likelihood\_combined} that the sampler calls.

\begin{figure}[htbp]
  \centering
  \includegraphics[width=\linewidth]{../gamma_scan_outputs/run_2026-05-24T10-19-32/plateau/plateau_loglike.png}
  \caption{$\log L(\gamma)$ scan at the OLS-truth $(c_0,\alpha,\beta)$ for each
  negative $\gtrue$ simulation.  Red dashed: $\gtrue$; black dotted: empirical
  peak; grey dashed: $\gamma = -0.5$. The y-axis is clipped to $-8000$ to
  reveal the peak structure; outside this window the log-likelihood drops by
  $10^9$--$10^{19}$ units as some clusters' implied $c$ saturates at
  $\gamma \lesssim -1.3$.}
  \label{fig:loglike}
\end{figure}

The peak positions and the log-likelihood gap between $\gtrue$ and the empirical
peak are summarised in Table~\ref{tab:loglike}.

\begin{table}[htbp]
  \centering
  \begin{tabular}{rrr}
    \toprule
    $\gtrue$ & peak $\gamma$ & $\log L(\gtrue) - \log L_{\rm peak}$\\
    \midrule
    $-1.0$ & $-0.525$ & $\mathbf{-3851}$\\
    $-0.8$ & $-0.550$ & $-276$\\
    $-0.6$ & $-0.425$ & $-1410$\\
    $-0.4$ & $-0.500$ & $-3$\\
    $-0.2$ & $-0.525$ & $-6$\\
    \bottomrule
  \end{tabular}
  \caption{Profile log-likelihood peak vs.\ $\gtrue$.  At $\gtrue = -1$ the
  likelihood prefers $\gamma \approx -0.5$ by $3851$ log-units. This is far
  larger than any sampler-level effect could explain.}
  \label{tab:loglike}
\end{table}

\noindent The plateau therefore lives in the likelihood, not in nautilus's
exploration.  A separate numerical feature visible in
Fig.~\ref{fig:loglike} --- a catastrophic drop in $\log L$ for
$\gamma \lesssim -1.3$ --- is an unrelated edge effect where some clusters'
$c = c_0 (\dm/0.1)^\gamma$ blow up past the regime where the NFW projection is
well-conditioned.

% ============================================================
\section{The bias is not a parametrisation artefact}
% ============================================================

A natural worry is that the formula
\begin{equation}
  \log_{10} c \;=\; \log_{10} c_0
  \;+\; \alpha (\log_{10} M - 14)
  \;+\; \beta \log_{10}\!\frac{1+z}{1.35}
  \;+\; \gamma \log_{10}\!\frac{\dm}{\dmref}
  \label{eq:cmodel}
\end{equation}
is the wrong functional form.  Two independent tests rule this out.

% ------------------------------------------------------------
\subsection{$\ghat$ is invariant under DELTA\_M\_REF (pure reparametrisation)}
% ------------------------------------------------------------

$\gamma$ is the slope w.r.t.\ $\log_{10}\dm$, so it is invariant under a rescaling
of $\dmref$ provided $c_0$ is rescaled by $(D_{\rm new}/D_{\rm old})^\gamma$.
\texttt{test\_delta\_m\_ref.py} verifies this on the $\gtrue=-1$ simulation:
the left panel of Fig.~\ref{fig:dmref} (consistent $c_0$, pure reparametrisation)
shows the $D=0.1$ and $D=0.236$ curves overlapping exactly with the same peak
at $\gamma = -0.55$.  The right panel keeps $c_0 = 4.575$ in both (mimicking the
tight $c_0 \sim N(4.57, 0.05)$ prior): at $D=0.1$ the peak is unchanged, but at
$D=0.236$ the prior-truth mismatch in the median-cluster $c$ shifts the peak to
$\gamma = -0.20$ --- a \emph{worse} bias.

\begin{figure}[htbp]
  \centering
  \includegraphics[width=\linewidth]{../gamma_scan_outputs/run_2026-05-24T10-19-32/plateau/delta_m_ref_test.png}
  \caption{Effect of $\dmref$ on the log-likelihood for $\gtrue = -1$.
  \emph{Left:} consistent $c_0(\gamma)$ --- the two curves are mathematically
  identical and both peak at $\gamma = -0.55$.
  \emph{Right:} $c_0$ held at $4.575$ (tight prior limit) --- the $D=0.236$ peak
  drifts to $\gamma = -0.20$, increasing the bias.}
  \label{fig:dmref}
\end{figure}

Shifting the $\dm$ pivot cannot move the plateau; if anything, choosing the
median makes it worse, because the tight $c_0$ prior was calibrated to the
$\dm=0.1$ pivot.

% ------------------------------------------------------------
\subsection{Even with the EXACT per-cluster truth concentrations, the model is
beaten by parametric $\gamma = -0.5$}
% ------------------------------------------------------------

This is the decisive test (\texttt{investigate\_truth\_vs\_data.py}).
For the $\gtrue=-1$ simulation we evaluate three log-likelihoods using the same
spherical NFW forward model, but supplying it with three different per-cluster
concentration arrays (Table~\ref{tab:truthvsparam}):

\begin{table}[htbp]
  \centering
  \begin{tabular}{lr}
    \toprule
    model & $\log L$\\
    \midrule
    spherical NFW with the \emph{exact per-cluster $c_{\rm truth}$} (no parametric $c$ relation) & $-487\,222.5$\\
    parametric model at $\gamma = \gtrue = -1$ with OLS $c_0,\alpha,\beta$ & $-487\,208.6$\\
    parametric model at $\gamma = \mathbf{-0.5}$ with the same OLS $c_0,\alpha,\beta$ & $\mathbf{-483\,357.7}$\\
    \bottomrule
  \end{tabular}
  \caption{Log-likelihoods for three $c$-assignment schemes on the same
  spherical NFW forward model. The parametric formula reproduces the truth-$c$
  likelihood to within $14$ log-units, but the same model with $\gamma = -0.5$
  beats them both by $\sim 3865$ log-units.}
  \label{tab:truthvsparam}
\end{table}

Two consequences follow:

\begin{enumerate}
\item \textbf{The parametric formula is not the bottleneck.}  Equation~\eqref{eq:cmodel}
  reproduces the truth-$c$ likelihood within $14$ log-units --- exactly as
  expected, since the OLS fit returns $\gamma_{\rm OLS} = -1.0000$ with
  $c_0 = 4.575$ and $\alpha,\beta$ at the canonical Dutton--Maccio values.
\item \textbf{The spherical NFW model with the truth $c$ values is still $3865$
  log-units worse than the same model at $\gamma = -0.5$.}  The likelihood is
  not asking for a different $c(M,z,\dm)$ law --- it is asking for a
  \emph{different per-cluster $c$ distribution}, one closer to canonical
  Dutton--Maccio values across the board.  No reparametrisation of $c$ can
  deliver this.
\end{enumerate}

\begin{figure}[htbp]
  \centering
  \includegraphics[width=\linewidth]{../gamma_scan_outputs/run_2026-05-24T10-19-32/plateau/truth_vs_param_profile.png}
  \caption{Mean $\Delta\Sigma$ profiles: data (triaxial, black) vs.\ spherical
  NFW with the per-cluster truth $c$ (purple), parametric model at $\gamma = \gtrue$
  (crimson), and parametric model at $\gamma = -0.5$ (steelblue).  At
  $\gtrue=-1$ the truth-$c$ spherical curve sits below the data, while the
  $\gamma=-0.5$ curve is closer --- that gap is the $3865$ log-unit preference.}
  \label{fig:truthvsparam}
\end{figure}

\begin{figure}[htbp]
  \centering
  \includegraphics[width=\linewidth]{../gamma_scan_outputs/run_2026-05-24T10-19-32/plateau/plateau_c_distribution.png}
  \caption{Model-implied $c_{200}$ distributions at three $\gamma$ values for
  each $\gtrue$ panel. At $\gtrue=-1$ the model histogram at $\gamma = \gtrue$
  (crimson) matches the truth (dotted black); at $\gamma=-0.5$ (steelblue)
  the distribution is much narrower, concentrated in $c \in [1, 10]$. The
  likelihood prefers the narrower distribution.}
  \label{fig:cdist}
\end{figure}

% ============================================================
\section{What the simulation actually generates: triaxial shape distribution}
% ============================================================

The lensing signal is computed with
\texttt{xlensing.sims.apply\_triaxial\_NFW\_shear\_region} using shapes drawn from
\texttt{sample\_triaxial\_shapes} with
\[
  b/a \sim U(0.5,\,1.0), \qquad
  c/a \sim U(0.3,\,b/a), \qquad
  \text{isotropic LOS}.
\]
Re-sampling the shapes for the $\gtrue=-1$ simulation (deterministic seed
$=10$) gives the distributions in Fig.~\ref{fig:triax} and the summary in
Table~\ref{tab:shapes}.

\begin{table}[htbp]
  \centering
  \begin{tabular}{lrrr}
    \toprule
    quantity & 16\,\% & median & 84\,\%\\
    \midrule
    $b/a$ & $0.582$ & $0.751$ & $0.921$\\
    $c/a$ & $0.364$ & $0.499$ & $0.692$\\
    triaxiality $T = (a^2-b^2)/(a^2-c^2)$ & $0.262$ & $\mathbf{0.643}$ & $0.869$\\
    projected $q_{\min}/q_{\max}$ & $0.512$ & $0.683$ & $0.860$\\
    \bottomrule
  \end{tabular}
  \caption{Triaxial shape distribution used in the simulation.}
  \label{tab:shapes}
\end{table}

The Franx-+91 triaxiality parameter $T$ classifies a halo as
\emph{oblate} ($T < 1/3$), \emph{generally triaxial} ($1/3 \le T \le 2/3$),
or \emph{prolate} ($T > 2/3$).  The simulated population partitions as

\begin{itemize}
  \item \textbf{$47\,\%$ prolate}
  \item $31\,\%$ generally triaxial
  \item $21\,\%$ oblate
\end{itemize}

\noindent The population is \emph{prolate-dominated} (median $T = 0.643$).
This is a direct consequence of the sampling rule: with $c/a \sim U(0.3, b/a)$
and $b/a \le 1$, the constraint $c \le b$ pushes $c/a$ low whenever $b/a$ is
low, producing a population where the \emph{minor} axis tends to be
significantly smaller than the major and intermediate axes --- geometrically,
this is prolate behaviour.  Only the corner of the $(b/a, c/a)$ plane near
$(1, 1)$ gives oblate shapes, and that corner has small measure.  Projected
on the sky, the median axis ratio is $q_{\min}/q_{\max} \approx 0.68$; there
is no perfectly circular projection in the sample
($q_{\min}/q_{\max}\to 1$ requires $c/a\to 1$, excluded by the prior).

\begin{figure}[htbp]
  \centering
  \includegraphics[width=\linewidth]{../gamma_scan_outputs/run_2026-05-24T10-19-32/plateau/triaxiality_distribution.png}
  \caption{Triaxial-shape distribution used in \texttt{gamma\_scan.py}.
  Top-left: intrinsic $b/a$ and $c/a$ marginals.  Top-right: triaxiality $T$
  histogram with the oblate / triaxial / prolate fractions labelled.
  Bottom-left: joint $(b/a, c/a)$ density with the prolate limit $c=b$
  highlighted.  Bottom-right: sky-projected axis ratio $q_{\min}/q_{\max}$.}
  \label{fig:triax}
\end{figure}

Stacked along an isotropic LOS, prolate halos with the major axis tilted toward
the observer produce a more centrally concentrated \emph{projected} $\Delta\Sigma$
than a spherical halo of the same intrinsic $c$; prolate halos perpendicular
to the LOS produce a less concentrated one.  The orientation-averaged signal
therefore differs from the spherical NFW signal at the same intrinsic $c$ ---
and the difference is precisely what the spherical inference model has to
absorb.

% ============================================================
\section{Why the mismatch produces a plateau at $\gamma \approx -0.5$}
% ============================================================

Three facts combine:

\textbf{(a)~The inference uses a spherical NFW for every cluster.}
\texttt{forward\_model\_combined} (in \texttt{destacking/model.py}) calls
\texttt{xlensing.model.NFW\_Delta\_Sigma}, which assumes spherical symmetry and
a single 3D concentration $c$. There is no orientation, no axis ratio, no
nuisance parameter that could absorb the projection of a triaxial halo.

\textbf{(b)~For prolate halos viewed along varying LOS, the projected
$\Delta\Sigma$ in a fixed radial range is not the same as a single spherical
NFW of the same intrinsic $c$.}  Empirically
(Fig.~\ref{fig:truthvsparam}), the orientation-averaged triaxial mean profile
at $\gtrue=-1$ is consistently \emph{above} the spherical NFW at the truth
$c$ values.  The data prefers a higher-amplitude inner profile at the same
fitted mass, which the spherical model can only reproduce by \emph{raising} the
effective $c$.

\textbf{(c)~Negative $\gtrue$ drives the $c$-distribution into bimodal tails.}
With $\gtrue=-1$ and $\dm$ spanning $[0.01, 2]$\,mag (median $0.236$):
\begin{itemize}
\item clusters with $\dm \gg 0.1$ (most of them; median $\dm/0.1 \approx 2.4$) get
$c$ suppressed below the Dutton--Maccio baseline, with many landing at
$c \in [0.5, 2]$;
\item clusters with $\dm \ll 0.1$ (about $10\,\%$ of them) get $c$ boosted above
$20$.
\end{itemize}
Both tails are regimes where the spherical-NFW projection is a poor match for
the orientation-averaged triaxial signal.  Very low $c$ gives nearly flat 3D
profiles whose 2D projection is dominated by the outskirts, where triaxial
shape effects are largest; very high $c$ gives nearly point-mass profiles whose
2D projection is dominated by the inner region, where prolate alignment causes
the biggest projection enhancement.

At $\gamma \approx -0.5$, by contrast, the same per-cluster $c$ distribution
stays in the well-behaved range $c \in [0.7, 8]$ for the vast majority of
clusters (steelblue histograms in Fig.~\ref{fig:cdist}).  In that regime, the
spherical-NFW model with appropriately tuned $c$ values approximates the
orientation-averaged triaxial signal much more closely.  The likelihood is just
measuring this: the spherical model fits the triaxial-averaged data best when
the $c$-distribution corresponds to $\gamma \approx -0.5$,
\emph{regardless of what $\gtrue$ actually is}.

This fixes the plateau location:
\begin{itemize}
\item Moving more negative than $\gamma \approx -0.5$ pushes more clusters into
the extreme-$c$ tails where the model-misspecification penalty grows
super-linearly (the catastrophic $\log L$ drop at $\gamma \lesssim -1.3$ is the
extreme limit of this).
\item Moving less negative gives $c$ values closer to canonical Dutton--Maccio
(better fit per cluster) but under-uses the $\dm$ correlation present in the
data.
\end{itemize}
The minimum of the model-misspecification penalty in the $\gamma$-direction
lands near $\gamma \approx -0.5$ for the current cluster sample.  The symmetric
plateau at $\ghat \approx +0.5$ on the positive side has the same mechanism with
the $\dm$ dependence flipped: positive $\gamma$ pushes low-$\dm$ clusters into
the low-$c$ tail.

% ============================================================
\section{Why we cannot fix this in the inference}
% ============================================================

The bias is an instance of an \emph{unmodelled nuisance}.  The simulation knows
each cluster's 3D shape and orientation; the inference does not, and the real
data does not either.  No observable in the survey --- single-cluster
$\Delta\Sigma$, richness $\lambda$, or $\dm$ --- recovers a cluster's individual
3D shape at the noise level we have.  Stacking would average out the
orientation but is exactly what the destacking pipeline is designed to avoid.

\textbf{No parametrisation of $c(M,z,\dm)$ can fix what is fundamentally a
model-geometry mismatch.}  In particular, none of the following address the
plateau:

\begin{itemize}
  \item changing \texttt{DELTA\_M\_REF} (verified empirically in \S 3.1);
  \item enlarging the $c_0$ prior (would let the sampler reach the
        spherical-equivalent $c$ distribution, but still locks $\ghat$ at the
        spherical-equivalent $\gamma \approx -0.5$);
  \item miscentering nuisance (already in the forward model; not the dominant
        effect);
  \item coherent covariance inflation (\texttt{coh\_amplitude}, since removed
        from the code base after producing only modest improvements; the bias
        survives any reasonable inflation amplitude).
\end{itemize}

What would in principle help but is outside the scope of this project, because
the real data does not constrain it:

\begin{itemize}
  \item triaxial NFW inference at the per-cluster level (no per-cluster
        orientation observable);
  \item a population-level triaxiality parameter marginalised over orientations
        (possible in principle; requires a robust prior, which the real data
        the pipeline is being designed for also does not constrain);
  \item projection-corrected $\Delta\Sigma$ templates (require knowing the
        dominant LOS-axis correlation in the survey, which is not measured).
\end{itemize}

% ============================================================
\section{Practical implication}
% ============================================================

When this pipeline is applied to real data, the recovered $\ghat$ will be the
value that fits the orientation-averaged signal under a spherical NFW
assumption, not the underlying intrinsic $\gtrue$.  To the extent that real
clusters share the prolate-dominated shape distribution of the simulation, the
same plateau-like bias will appear.  The right way to report the result is
therefore as the \emph{projection-averaged effective transfer function}
$\gamma_{\rm eff}(\gtrue)$ --- exactly the curve the gamma scan produces ---
rather than treating $\ghat$ as an unbiased estimator of intrinsic $\gtrue$.

% ============================================================
\section*{Files}
% ============================================================

\begin{itemize}
  \item Reference run: \texttt{gamma\_scan\_outputs/run\_2026-05-24T10-19-32/}
  \item Plateau plots (under \texttt{plateau/}):
    \texttt{plateau\_loglike.png},
    \texttt{plateau\_c\_distribution.png},
    \texttt{plateau\_mean\_profile.png},
    \texttt{truth\_vs\_param\_profile.png},
    \texttt{delta\_m\_ref\_test.png},
    \texttt{triaxiality\_distribution.png}.
  \item Diagnostic scripts (under \texttt{scripts/}):
    \texttt{investigate\_gamma\_plateau.py},
    \texttt{investigate\_truth\_vs\_data.py},
    \texttt{test\_delta\_m\_ref.py},
    \texttt{plot\_triaxiality\_distribution.py}.
  \item Transfer-function summary plots:
    \texttt{scan\_transfer\_function.png},
    \texttt{scan\_combined\_posteriors.png}.
\end{itemize}

\end{document}
